\section{Background}
\label{sec:background}

Pedestrian evacuation models have a long physical tradition. The social-force
model represents pedestrians as self-driven particles subject to destination,
repulsion, and contact forces~\cite{helbing1995social}. Cellular automata
models discretize space and update local movement decisions through transition
rules~\cite{burstedde2001simulation}. Evacuation-model reviews emphasize that
these physical models must be interpreted together with scenario assumptions,
calibration data, and validation scope rather than as universally predictive
engines~\cite{gwynne1999review,schadschneider2009evacuation}.
Fundamental-diagram work relates crowd density to walking speed and flow
capacity~\cite{fruin1971pedestrian,weidmann1993transporttechnik}.
These tools are useful because evacuation failure is often physical:
bottlenecks, counter-flow, density waves, and slower-is-faster effects can
dominate individual intent.

Hazard-coupled evacuation adds another layer. In CBRN or smoke scenarios, the
environment changes while people move through it; visibility, speed, and
physiology can degrade as a function of exposure. Prior toxic-gas evacuation
models already show that information perception and transmission matter for
route choice under a spreading contaminant~\cite{makmul2020cellular,liu2021toxicgas}.
Chiyoda builds on this insight but makes the communication policy itself the
experimental object.

Information-aware evacuation models also intersect with route-choice theory,
game theory, and bounded rationality. Pedestrian route-choice models connect
activity constraints, perceived costs, and network structure~\cite{hoogendoorn2004routechoice},
while evacuation experiments show that visibility can change route-selection
behavior~\cite{guo2012visibility}. Bayesian formulations can model incomplete information
and route-choice equilibria under congestion~\cite{mah90dynamic,wang2023bayesian}. Fire
behavior and virtual-evacuation studies also show that exit choice is shaped
by affiliation, social influence, stress, familiarity, and other occupants'
movement~\cite{sime1983affiliative,kobes2010building,bode2013human,kinateder2014social}. Social
group ABMs review the importance of communication, affiliation, and collective
behavior, while also noting persistent gaps in behavioral realism and
validation~\cite{templeton2023social,senanayake2024agent}. Chiyoda takes a
simulation-first position: it does not claim to solve empirical validation in
one paper, but it makes belief uncertainty measurable and interventions
reproducible.

The information-theoretic lens uses Shannon entropy~\cite{shannon1948mathematical}.
For each agent, Chiyoda computes entropy over exit existence, hazard knowledge,
and general danger. Global entropy is the population mean. Entropy is not a
normative objective by itself: the goal is to test when reducing entropy
improves safety and when it merely concentrates the crowd's behavior.
